% ============================================================================
% DISSENSUS PREPRINT TEMPLATE v3.0.0
% ============================================================================
% Springer Nature-inspired academic preprint with minimal personal branding.
% Derived from ROM paper (arXiv:2601.06363) — the proven, compiled reference.
%
% Author: Murad Farzulla | Dissensus
% Last Updated: February 2026
%
% PAPER: Privacy-Preserving Financial Surveillance
% ============================================================================

\documentclass[11pt]{article}

% ============================================================================
% PACKAGES
% ============================================================================

% Typography and layout
\usepackage[T1]{fontenc}
\usepackage{mathptmx}                    % Times (most universal academic font)
\usepackage[a4paper, margin=1in]{geometry}
\usepackage{setspace}

% Graphics and figures
\usepackage{graphicx}
\usepackage{float}

% Tables
\usepackage{booktabs}
\usepackage{array}
\usepackage{multirow}
\usepackage{longtable}

% Mathematics
\usepackage{amsmath}
\usepackage{amssymb}
\usepackage{amsthm}
\renewcommand{\qedsymbol}{$\blacksquare$}  % Filled QED squares (Springer style)

% Bibliography
\usepackage[round,authoryear]{natbib}
\bibliographystyle{plainnat}

% Colors — burgundy is the only branding touch
\usepackage{xcolor}
\definecolor{brandburgundy}{RGB}{128,0,32}

% Hyperlinks and PDF metadata
\usepackage{url}
\usepackage[colorlinks=true,
            linkcolor=brandburgundy,
            citecolor=brandburgundy,
            urlcolor=brandburgundy,
            breaklinks=true,
            pdftitle={What Does Identity Access Buy? Pricing the Surveillance Increments in CBDC Transaction Monitoring},
            pdfauthor={Murad Farzulla and Andrew Maksakov},
            pdfkeywords={CBDC, privacy, AML, identity, detection, measurement, zero-knowledge proofs}]{hyperref}

% Allow URLs to break at hyphens and slashes
\def\UrlBreaks{\do\/\do-\do_}
\expandafter\def\expandafter\UrlBreaks\expandafter{\UrlBreaks\do\a\do\b\do\c\do\d\do\e\do\f\do\g\do\h\do\i\do\j\do\k\do\l\do\m\do\n\do\o\do\p\do\q\do\r\do\s\do\t\do\u\do\v\do\w\do\x\do\y\do\z}

% Formatting
\usepackage{enumitem}
\usepackage{fancyhdr}

% Section formatting — burgundy headings with proper spacing
\usepackage{titlesec}
\titleformat{\section}{\normalfont\large\bfseries\color{brandburgundy}}{\thesection}{0.5em}{}
\titleformat{\subsection}{\normalfont\normalsize\bfseries\color{brandburgundy}}{\thesubsection}{0.5em}{}
\titleformat{\subsubsection}{\normalfont\small\bfseries\color{brandburgundy}}{\thesubsubsection}{0.5em}{}
\titlespacing*{\section}{0pt}{2ex plus 0.8ex minus 0.2ex}{1ex plus 0.3ex}
\titlespacing*{\subsection}{0pt}{1.5ex plus 0.5ex minus 0.2ex}{0.8ex plus 0.2ex}
\titlespacing*{\subsubsection}{0pt}{1.2ex plus 0.4ex minus 0.2ex}{0.6ex plus 0.2ex}

% Headers/footers — minimal (page number only)
\pagestyle{fancy}
\fancyhf{}
\fancyfoot[C]{\small\thepage}
\renewcommand{\headrulewidth}{0pt}
\renewcommand{\footrulewidth}{0pt}

\fancypagestyle{firstpage}{
  \fancyhf{}
  \fancyfoot[C]{\small\thepage}
  \renewcommand{\headrulewidth}{0pt}
  \renewcommand{\footrulewidth}{0pt}
}

% ============================================================================
% PAPER METADATA
% ============================================================================

\newcommand{\papernum}{DAI-2511}              % Working paper number
\newcommand{\paperver}{2.0.0}
\newcommand{\paperdate}{February 2026}
\newcommand{\paperdoi}{10.5281/zenodo.17917938}
\newcommand{\paperurl}{https://systems.ac/2/DAI-2511}  % ASCRI page

% ============================================================================
% DOCUMENT BODY
% ============================================================================

\begin{document}
\setstretch{1.15}
\thispagestyle{firstpage}

\begin{center}
{\small\textsc{Dissensus Working Paper Series}}\\[0.2em]
{\small \href{\paperurl}{\papernum}}
\end{center}

\vspace{1.5em}

\begin{center}
{\LARGE\bfseries What Does Identity Access Buy?\\Pricing the Surveillance Increments in CBDC Transaction Monitoring}\\[1.5em]

{\large Murad Farzulla}\textsuperscript{1,2,*}, \quad {\large Andrew Maksakov}\textsuperscript{1}\\[0.8em]

{\small
  \textsuperscript{1}\href{https://dissensus.ai}{Dissensus}, London, UK \quad
  \textsuperscript{2}King's College London, London, UK%
}\\[0.5em]

{\footnotesize
  \textsuperscript{*}Correspondence: \href{mailto:murad@dissensus.ai}{murad@dissensus.ai}
  \quad
  ORCID: \href{https://orcid.org/0009-0002-7164-8704}{0009-0002-7164-8704}%
}\\[0.3em]
{\footnotesize \paperdate}
\end{center}

\begin{abstract}
\noindent Proposals for retail central bank digital currency increasingly accept that
anti-money-laundering obligations require identity-linked transaction monitoring. The cost of that
requirement has not been measured. Central-bank experiments that address privacy-preserving
detection---the BIS Innovation Hub's Project Aurora, and Hertha with the Bank of England---ablate the
\emph{data-sharing} axis: what is gained when institutions pool transactions. Neither ablates the
\emph{identity} axis: what is gained from knowing who owns an account, over and above knowing which
accounts move together. This paper measures that increment. We build a tiered ablation harness in
which four nested evidence tiers---structure only on unlinked pseudonyms, plus pseudonymous linkage,
plus identity attributes, plus watchlist access---are scored on identical data with entity-disjoint
folds and entity-clustered bootstrap intervals. A pre-specified degeneracy audit fails the run if any
single feature separates the classes at entity-level AUC above 0.95; the disjoint wallet-count ranges
of an earlier version of this work put wallet count at AUC 1.000 by construction, which the audit
rejects. Equivalence between
tiers is tested by TOST rather than inferred from a null result, and a falsification world in which
surveillance is built to win must return that verdict or the pipeline aborts. On synthetic data the
answer depends on the detector: for a gradient-boosted model, pseudonymous linkage carries the
increment and full identity access is statistically equivalent to linkage alone, robustly across a
swept generating parameter; for a logistic model, identity access is decisively superior on average
precision at every point of that sweep. In the regimes tested, then, the marginal value of identity
access is model- and metric-dependent rather than a fixed property of the data. We report the
bounds this evidence cannot cross---synthetic generation, a few hundred positive cases, two model
families, and the fact that no identified public anti-money-laundering benchmark jointly exposes the
behavioural, linkage, identity and watchlist axes that the full ablation requires, so the quantity
cannot presently be validated on real data---and set out the architecture and governance this
measurement motivates and constrains, and those it does not.

\vspace{0.5em}
\noindent\textbf{Keywords:} central bank digital currency, privacy, financial surveillance,
anti-money laundering, identity, zero-knowledge proofs, detection, measurement

\vspace{0.3em}
\noindent\textbf{JEL Codes:} E42, E58, G28, K42, O33
\end{abstract}

\vspace{1em}

\section*{Research Context}

This work forms part of the Adversarial Systems Research program, which investigates stability,
alignment, and friction dynamics in complex systems where competing interests generate structural
conflict. The privacy-preserving CBDC question addressed here is one such dynamic: the tension between
state surveillance capacity and citizen financial autonomy. The paper extends prior work on AML
regulatory failure \citep{farzulla2025hedging}, which documented how current frameworks target
primitive laundering while sophisticated actors exploit derivatives and offshore structures with
impunity; where that analysis diagnosed the problem, this paper asks what one of its assumed remedies,
identity-linked monitoring, is actually worth.

\section*{Acknowledgments}

The authors acknowledge David Chaum's foundational work on eCash and blind signatures, and the
subsequent CBDC proposal with Grothoff and Moser \citep{ChaumGrothoffMoser2021} demonstrating that
privacy-preserving central bank money is institutionally viable. Andrew Maksakov contributed the
PET--AML stack simulation (\texttt{pet\_aml\_sim.py}, Section~\ref{sec:feasibility}). This paper
benefited from collaboration with Claude (Anthropic) on analysis, literature synthesis, and drafting;
the authors take full responsibility for all claims, errors, and interpretive choices. Related work in
the series is available through the Adversarial Systems \& Complexity Research Initiative
(\href{https://systems.ac}{ASCRI}). Correspondence:
\href{mailto:murad@dissensus.ai}{murad@dissensus.ai}.

\clearpage
\tableofcontents
\thispagestyle{firstpage}
\clearpage
\section{Introduction}
\label{sec:introduction}

Every serious proposal for a retail central bank digital currency accepts, more or less explicitly,
that anti-money-laundering obligations require identity-linked transaction monitoring. The design
question is then taken to be how much privacy can be preserved \emph{around} that requirement---how
to encrypt, tier, threshold, or delay access to identity data whose collection is treated as
settled. The requirement itself is rarely examined, and its price has never been measured. This
paper measures it.

The requirement is worth examining because the surveillance it justifies performs poorly on its own
terms. Estimates of laundered funds sit at two to five percent of global output \citep{UNODC2011},
and of that flow the fraction interdicted is very small: reviewing the evidence, \citet{Pol2020}
concludes that seizures run at ``some undetermined fraction of one percent,'' a rate at which, in his
words, ``the disruption of criminal funds hardly constitutes a rounding error.'' A system that
imposes identity monitoring on an entire population to recover under one percent of the target flow
is not obviously buying much with the identities it collects, and the burden should be on the
architecture to show that it is.

The enforcement record compounds the doubt, because the friction and the failure fall on different
people. Structuring, money mules, and basic layering---the primitive methods available to
individuals, often coerced ones---are prosecuted vigorously, while the institutions through which the
largest flows pass are resolved with financial penalties. When Danske Bank's Estonian branch was
found to have processed on the order of \$160--200 billion in suspicious non-resident transactions,
the resolution was a guilty plea and more than \$2 billion in forfeiture to the U.S. Department of
Justice, with a further \$413 million to the SEC \citep{Lynch2022}; the individuals criminally
charged were branch-level staff rather than group leadership. Two billion dollars is a large number
and roughly a cent on each dollar of the flow it settled. The asymmetry is not that enforcement is
absent but that it lands hardest on the least sophisticated participants, which is precisely the
population for whom identity monitoring is most totalizing and least necessary.

Public resistance to that monitoring is well documented and specific. Respondents to the European
Central Bank's digital-euro consultation ranked privacy the most important feature of the instrument,
ahead of security \citep{ECBConsultation2021}; in a Cato Institute/YouGov survey, 74 percent of U.S. adults said
they would oppose a CBDC if it let the government control how they spend, and 68 percent if it let the
government monitor their spending \citep{CatoCBDC2023}. Whatever one makes of advocacy-sponsored
polling, the direction is consistent across sources and the objection is not to digital money as such
but to the identity-linked visibility its proposed architectures assume.

So the question this paper asks is narrow and, we think, prior to the architecture debate: granted a
detector operating on transaction data, what is the marginal value of each additional increment of
identity access? Concretely, we separate three increments a monitoring regime might be granted in
sequence---the ability to resolve pseudonymous wallets to a common owner (\emph{linkage}), the
attachment of verified identity attributes to that owner (\emph{identity}), and consultation of a
watchlist (\emph{watchlist})---and measure what each buys in detection performance over the one below
it, holding the detector and the data fixed. The decomposition is the contribution: it isolates the
value of \emph{identity per se}, above and beyond the linkage of activity, which no prior evaluation
we are aware of has done (Section~\ref{sec:already-known}).

The answer, developed in Sections~\ref{sec:measurement-problem}--\ref{sec:results}, is that the
marginal value of identity access is not a property of the data but of the detector applied to it. A
detector capable of exploiting the interaction structure of behavior recovers almost everything from
pseudonymous linkage alone, and identity attributes add nothing measurable on top; a
capacity-limited detector cannot recover that structure and leans on identity attributes to
substitute for what it cannot infer. Detection capability and identity collection are, to a
meaningful degree, substitutes. We establish this on a purpose-built measurement instrument rather
than by architectural argument, we show it is robust across the range of a swept generating
parameter rather than an artifact of one synthetic world, and we are explicit about the ceiling on
the evidence: no public anti-money-laundering dataset carries an identity axis, so the quantity we
decompose cannot at present be measured on real data by anyone, which is exactly why a controlled
synthetic instrument is required to pose the question at all.

The remainder of the paper is organized around that instrument. Section~\ref{sec:already-known}
places the work against the CBDC design literature and the cryptographic-compliance lineage it builds
on. Section~\ref{sec:measurement-problem} states why a naive ablation manufactures its own answer,
using a failed earlier version of this work as the cautionary case.
Section~\ref{sec:instrument} describes the instrument and its guards, and
Section~\ref{sec:results} reports what it finds, including the robustness surface and the
falsification check. Sections~\ref{sec:licenses}--\ref{sec:feasibility} then set out the
architecture and cryptographic stack this measurement licenses---and the parts of it the measurement
does \emph{not} support---followed by governance (Section~\ref{sec:governance}) and the limitations
that bound the whole (Section~\ref{sec:limitations}).


\section{What Is Already Known}
\label{sec:already-known}

\subsection{CBDC design and the place of privacy}

It is tempting to motivate a privacy-preserving architecture by asserting that the CBDC design
literature treats surveillance as a default and privacy as an afterthought. That assertion does not
survive contact with the literature, and we do not make it. \citet{AuerBohme2020}, whose taxonomy of
direct, indirect, and hybrid architectures is the field's common reference, make central-bank
visibility the very axis that distinguishes their categories rather than a constant across them: in
the indirect model the central bank ``keeps no record of individual claims,'' while only in the
direct model does it hold ``a record of all balances.'' They treat privacy as a design variable to
be engineered, proposing that users ``retain cryptographic proofs of their CBDC balances, rather than
oblige the central bank to hold them,'' and endorsing ``unlinkable pseudonyms, as envisaged in
Chaum's pioneering work''---which is to say, the architecture this paper later specifies. The
comprehensive review by \citet{AuerFrostGambacorta2022} likewise organizes itself partly around
privacy, ties the privacy question directly to architecture, and observes that a design ``in which
the central bank has no information on retail transactions'' is available. The design literature has
not ignored privacy; it has, if anything, anticipated the architectural direction taken here.

What the design literature does tend to do, and what the framing documents make explicit, is treat
the privacy--surveillance relation as a balance to be struck rather than a quantity to be measured.
The ECB frames the goal as the ``appropriate balance between privacy and other public policy
objectives, like countering money laundering'' \citep{ECB2024Blog}; the Federal Reserve's discussion
paper poses the design problem as striking ``an appropriate balance $\ldots$ between safeguarding the
privacy rights of consumers and affording the transparency necessary to deter criminal activity''
\citep{FedReserve2022}. The metaphor of balance presupposes that the two quantities trade off along a
known frontier. Whether they do, and at what exchange rate, is an empirical question that the framing
documents pose but do not answer---and that this paper takes up.

\subsection{Privacy-preserving compliance is a mature research program}

Nor is the idea of reconciling compliance with privacy new. A cryptographic-compliance lineage
stretches back three decades, from the trustee-based e-cash of \citet{CamenischMaurerStadler1996},
through compact e-cash \citep{Camenisch2005CompactEcash} and accountable privacy for decentralized
ledgers \citep{Garman2016Accountable}, to recent systems that combine regulatory tracing with
transaction privacy---Platypus \citep{Wuest2022Platypus}, PEReDi \citep{Kiayias2022PEReDi}, and the
Privacy Pools construction \citep{Buterin2024PrivacyPools}. Each of these establishes that selective,
warranted, or threshold-gated disclosure can coexist with default transaction privacy. The
contribution of the present work is not to add another such construction; the primitives it later
assembles (Section~\ref{sec:feasibility}) are drawn from this literature rather than invented here.

\subsection{The unmeasured axis}

The gap this paper addresses is narrower and empirical. Central banks have begun to measure
privacy-preserving detection directly, at deployment scale: the BIS Innovation Hub's Project Aurora
\citep{BISAurora2023} and the Hertha collaboration with the Bank of England \citep{BISHertha2025}
both quantify what privacy-preserving analytics recover relative to full-data access. But both ablate
the same axis---\emph{data sharing}, the gain from pooling observations across institutions---and
neither separates the \emph{identity} axis from the \emph{linkage} axis beneath it. Neither asks what
knowing who owns an account buys over and above knowing which accounts move together. The nearest
work in the academic literature that varies identity information, which ablates socio-demographic
profile attributes for transaction classification, does so on an already entity-resolved graph with
no unlinked-pseudonym tier and no watchlist axis, so it cannot make the separation this paper
targets. Isolating the value of identity per se---distinct from linkage below it and watchlist access
above it---is, to our knowledge, unmeasured, and the rest of this paper measures it.
% =====================================================================
% CBDC paper v4 — NEW SPINE (Sections 3, 4, 5 of the restructured paper)
% Written 24 Jul 2026 against detection/results/ (seed 20260707).
% Every number in Section 5 is read directly from results_*.csv,
% tost_*.json and degeneracy_audit_*.json. Nothing here is asserted.
% US spelling per house rule. natbib \citep/\citet, booktabs.
% =====================================================================

\section{The Measurement Problem}
\label{sec:measurement-problem}

The quantity we want is the marginal value of each surveillance increment: how much detection
performance is gained by resolving pseudonymous wallets to entities, how much is gained on top of
that by attaching verified identity attributes to those entities, and how much is gained again by
consulting a watchlist. Stated that way the measurement looks straightforward---train a detector on
nested feature sets and difference the results. It is not, and the reason it is not is the central
methodological problem of this paper.

An ablation over synthetic data measures the data-generating process, not the world. If the process
that labels an entity a launderer also endows it with a distinguishing marginal feature available at
some tier, that tier wins by construction, and the ablation reports the analyst's modeling choice
back to them with a confidence interval attached. The failure is not subtle, but it is easy to
commit and hard to see from inside, because the resulting numbers look like findings.

We can be concrete, because an earlier version of this work committed exactly this error. In that
design, laundering entities were specified to control three to six wallets and legitimate entities
one to two, with no overlap between the ranges. Wallet count therefore separated the two classes
perfectly, and any detector with access to it achieved an area under the ROC curve of $1.000$. The
reported finding---that watchlist access added no marginal value, moving detection from $1.00$ to
$1.00$---was arithmetically forced rather than observed: a saturated ceiling admits no improvement,
so the ablation could not have returned anything else. The finding was not weak evidence for the
paper's thesis. It was no evidence at all, and it pointed in the thesis's favor.

The error generalizes in both directions, which is what makes it dangerous rather than merely
embarrassing. A process that places the discriminating signal in a structural feature manufactures
the conclusion that surveillance is unnecessary; a process that places it in an identity attribute
manufactures the conclusion that surveillance is essential. Neither conclusion survives contact
with a differently specified process, and an ablation reported without a test for this condition
carries no information about which case obtains.

An instrument adequate to the question therefore needs four properties, and the rest of this
section states them as requirements before Section~\ref{sec:instrument} describes how they are met.
First, laundering must be a \emph{latent behavior} realized over time---a sequence of placement,
layering, and integration decisions---rather than a labeled trait that some feature reads off
directly. Second, the instrument must \emph{test itself for degeneracy} before any headline number
is computed, and must fail loudly rather than warn. Third, the equivalence of two tiers must be
established by an \emph{equivalence test}, since failing to reject a difference is not evidence that
none exists, and with a small positive class it is very weak evidence indeed. Fourth, the instrument
must be shown to be \emph{capable of returning the opposite answer}: a harness that reports
``surveillance adds nothing'' on every input it is given is not measuring anything.


\section{Instrument}
\label{sec:instrument}

\subsection{Generating process}

Entities transact over a simulated 120-day window. Of 8{,}000 entities, a base rate of $0.05$
designates launderers; 388 are realized in the seeded run reported here. Each entity controls a number of
wallets drawn from a \emph{single shared distribution} applied to both classes, which is the
minimal fix for the failure described above: wallet count carries no class information by
construction, and any predictive power it acquires must come from how wallets are used.

Laundering is realized as behavior rather than assigned as an attribute. A laundering entity runs
placement, layering, and integration cycles over the window, splitting value across its wallets,
routing through external counterparties, and structuring individual transfers beneath a reporting
threshold with probability $0.6$. An obfuscation parameter of $0.6$ governs how aggressively the
entity spreads activity to frustrate aggregation. Legitimate entities transact under the same
mechanics without the cycle structure. The classes are consequently distinguished by \emph{patterns
of use} that a detector must reconstruct, not by any quantity it can read.

Identity attributes are drawn from overlapping class-conditional distributions rather than disjoint
ones: KYC tier, account age, prior suspicious-activity-report count, and a jurisdiction risk score
each differ in distribution between classes while assigning no entity a determinative value. The
watchlist is imperfect in both directions, with a true-positive rate of $0.4$ and a false-positive
rate of $0.02$, which is the property that makes the fourth tier a realistic increment rather than
a label leak.

\subsection{Evidence tiers}

Four nested tiers correspond to four levels of access a monitoring regime might be granted:

\begin{itemize}
  \item \textbf{T1 --- structure only.} Sixteen per-wallet flow and timing features---volumes,
        degrees, inter-transfer lags, burstiness, proximity to the reporting threshold---modeled at
        \emph{wallet} level, with an entity scored by the maximum score across its wallets. The
        regime observes individual pseudonymous accounts and can flag one that behaves badly, but
        cannot group accounts by owner.
  \item \textbf{T2 --- plus pseudonymous linkage.} Modeling moves to entity level over nineteen
        features that only linkage makes computable: wallet count, cross-wallet means and maxima,
        internal transfer volume, chained fast-follow value, external counterparty count. The regime
        knows which accounts move together but not who owns them.
  \item \textbf{T3 --- plus identity attributes.} Verified KYC tier, account age, prior SAR count,
        and jurisdiction risk are attached to each entity.
  \item \textbf{T4 --- plus watchlist.} A watchlist membership bit is added.
\end{itemize}

The upper three tiers are strictly nested at entity level, $\mathrm{T2} \subset \mathrm{T3} \subset
\mathrm{T4}$, so the identity and watchlist increments are each the value of exactly one additional
access right holding everything else fixed. T1 is deliberately \emph{not} a subset of T2: it is a
different unit of analysis, because the absence of linkage is not the removal of a column but the
inability to form the entity-level row at all. The T1$\to$T2 step therefore prices a change of
modeling regime rather than a change of feature set, and we report it as such.

One consequence of this design should be stated plainly rather than left for a reader to find. Even
at T1, the \emph{evaluation} is entity-level: wallets inherit their owner's cross-validation fold,
which is required to prevent an entity's wallets from straddling the train and test partitions, and
entity scores are formed by aggregating across an entity's wallets. The T1 detector is denied
linkage as a feature; the T1 protocol still assumes ground-truth linkage for scoring. T1 therefore
measures a regime that cannot \emph{exploit} linkage rather than one in which linkage is
unavailable in principle, and the true structure-only floor is somewhat below the figure reported
here.

This is the decomposition the paper contributes: the BIS Innovation Hub's Project Aurora
\citep{BISAurora2023} and the Hertha collaboration with the Bank of England
\citep{BISHertha2025} both ablate the \emph{data-sharing} axis, asking what
is gained when institutions pool observations, and neither separates the \emph{identity} axis from
the linkage axis beneath it.

Two model classes are fit at every tier: a logistic regression and a histogram-based gradient
boosting classifier. This is not a horse race, and we do not report a winner. The contrast between a
model of limited capacity and one able to exploit interactions among behavioral features is itself
the finding of Section~\ref{sec:results}, and a single-model design would have concealed it.

\subsection{Evaluation and its guards}

Scoring is at entity level with entity-disjoint cross-validation folds: no wallet belonging to an
evaluated entity appears in training, which prevents the linkage tier from leaking its own answer.
Confidence intervals are entity-clustered bootstraps, so the interval reflects the number of
independent entities rather than the number of transactions. We run 8{,}000 entities at a base rate
of $0.05$, yielding on the order of 400 realized launderers; an earlier version of this work ran
800, and Section~\ref{sec:results} shows why that was too few---at $\sim$29 positive cases every
average-precision comparison is inconclusive, and the study's whole verdict rested on a single
saturating metric.

Performance is measured on two metrics, not one, and the distinction turns out to carry the result.
Area under the ROC curve saturates at these operating points---once the study is adequately powered
every tier sits above $0.98$ for the capable model, so ``equivalent within a small margin'' is close
to unfalsifiable and the verdict conveys little. Average precision does not saturate at $\sim$5\%
prevalence, and it is the quantity a monitoring regime actually operates on when it allocates a fixed
alert budget. The two can disagree, and where they do the disagreement is informative rather than a
nuisance (Section~\ref{sec:results}).

A \textbf{pre-specified degeneracy audit} runs before any headline number and terminates the run if
any single feature achieves entity-level AUC above $0.95$, on the reasoning that a lone feature that
nearly separates the classes indicates the generating process has leaked the label. In the runs
reported here the worst single feature reaches $0.840$ (mean fraction of fast-follow transfers) in
the primary world and $0.887$ (jurisdiction risk) in the falsification world, both below the
threshold. The disjoint wallet-count ranges of the earlier version described in
Section~\ref{sec:measurement-problem} put wallet count at an entity-level AUC of $1.000$ by
construction---perfect separation of two non-overlapping integer ranges---so the audit rejects that
generating process outright and refuses to produce results at all.

Tier equivalence is assessed by two one-sided tests. We fix the margin at $\delta = 0.03$; this was
chosen on the AUC scale before any average-precision result existed, so applying the same number to
average precision---whose baseline is prevalence rather than $0.5$---would be a post hoc margin
choice on a different scale, and we flag it as such wherever it appears. To avoid resting anything on
that choice, every comparison also reports its \emph{equivalence bound}: the smallest margin at which
the test would declare equivalence, which is margin-free and lets a reader apply their own threshold.
We report equivalence as a positive finding only where the test supports it, \textsc{inconclusive}
where it does not, and we never infer equivalence from a failure to reject a difference.

Two guards close the loop against the instrument flattering its author. First, the harness must
demonstrate that it \emph{can} return the opposite verdict: a second world,
\texttt{surveillance\_strong}, is parameterized so identity attributes genuinely dominate behavioral
structure, and the pipeline exits non-zero unless that world returns \textsc{surveillance\_superior}
on at least one comparison, so a regression that biased the instrument toward the paper's thesis
breaks the build rather than producing a publishable number. Second, because the deepest objection to
any synthetic ablation is that the analyst chose the generating process, the primary result is
reported not at one operating point but across a sweep of the process's most consequential free
parameter (Section~\ref{sec:robustness}), so what the paper claims is an invariant of the sweep
rather than a fact about a single world.

\section{Results}
\label{sec:results}

All numbers in this section come from the seeded pipeline at the canonical operating point (8{,}000
entities, base rate $0.05$, 388 realized launderers, seed 20260707), except
Section~\ref{sec:robustness}, which sweeps the generating process, and Section~\ref{sec:falsify},
which reports the falsification world. Every figure regenerates from committed code.

\subsection{The operating point}
\label{sec:operating-point}

Table~\ref{tab:tier-perf} reports detection performance by tier for both models on both metrics, with
entity-clustered 95\% bootstrap intervals. Reading the two metrics side by side is the point of the
table, not a formality: they tell different stories about the same scores.

\begin{table}[htbp]
\centering
\caption{Detection performance by evidence tier (8{,}000 entities, 388 launderers, seed 20260707).
Intervals are entity-clustered bootstraps. AUC saturates for the capable model; average precision
(AP) does not.}
\label{tab:tier-perf}
\begin{tabular}{llcc}
\toprule
Model & Tier & AUC (95\% CI) & AP (95\% CI) \\
\midrule
Logistic & T1 structure only & 0.965 [0.951, 0.977] & 0.851 [0.819, 0.880] \\
         & T2 + linkage      & 0.981 [0.975, 0.987] & 0.859 [0.829, 0.886] \\
         & T3 + identity     & 0.987 [0.983, 0.992] & 0.894 [0.868, 0.918] \\
         & T4 + watchlist    & 0.992 [0.989, 0.995] & 0.922 [0.899, 0.943] \\
\addlinespace
Gradient & T1 structure only & 0.987 [0.981, 0.992] & 0.912 [0.891, 0.933] \\
boosting & T2 + linkage      & 0.998 [0.996, 0.999] & 0.980 [0.971, 0.989] \\
         & T3 + identity     & 0.998 [0.997, 0.999] & 0.981 [0.971, 0.989] \\
         & T4 + watchlist    & 0.999 [0.998, 0.999] & 0.985 [0.977, 0.992] \\
\bottomrule
\end{tabular}
\end{table}

The AUC columns for the gradient-boosted model are visually flat---0.987 to 0.999 across four
tiers---and this is the trap the earlier version of this work fell into. When every tier scores above
0.98, differences among them are compressed into the last percent of a bounded scale, and any
equivalence test with a plausible margin will return equivalence almost regardless of what the tiers
actually contain. Average precision, which weights the ranking of the $\sim$5\% of entities that are
positive rather than the whole ROC surface, keeps its resolution: the logistic model moves from 0.851
to 0.922 across the tiers, room in which real differences can show.

\subsection{The increments, and the mechanism behind them}
\label{sec:increments}

Table~\ref{tab:increments} differences adjacent tiers to give the marginal value of each access right.
We read the average-precision column, for the reason just given.

\begin{table}[htbp]
\centering
\caption{Marginal gain from each additional access right at the canonical operating point. Read the AP
column; the AUC column is shown to make the saturation visible.}
\label{tab:increments}
\begin{tabular}{llccc}
\toprule
Model & Metric & Linkage (T1$\to$T2) & Identity (T2$\to$T3) & Watchlist (T3$\to$T4) \\
\midrule
Logistic & AP  & $+0.009$ & $+0.035$ & $+0.027$ \\
         & AUC & $+0.016$ & $+0.006$ & $+0.005$ \\
\addlinespace
Gradient boosting & AP  & $+0.069$ & $+0.000$ & $+0.004$ \\
                  & AUC & $+0.011$ & $+0.000$ & $+0.001$ \\
\bottomrule
\end{tabular}
\end{table}

The two models allocate the gain in opposite ways, and the contrast is the finding. On average
precision the gradient-boosted model takes $+0.069$ from linkage and then nothing from identity
($+0.000$) and almost nothing from the watchlist ($+0.004$). The logistic model takes almost nothing
from linkage ($+0.009$) and then $+0.035$ from identity and $+0.027$ from the watchlist.

This is not two noisy estimates of one quantity; it is a mechanism. The signal that separates
launderers from legitimate entities lives in the \emph{interaction structure} of behavior across an
entity's wallets---how fast value chains between accounts, how activity bursts and disperses. A
gradient-boosted model can exploit that structure once linkage lets it see all of an entity's wallets
together, which is why linkage carries its entire gain and identity attributes add nothing a
sufficiently capable model has not already recovered from behavior. A linear model cannot represent
those interactions, so linkage buys it little; it leans instead on identity attributes, which supply
the class signal directly rather than through structure it cannot model. Identity access substitutes
for a capability the weak model lacks and the strong model does not need.

\subsection{Equivalence, and where the metrics disagree}
\label{sec:equivalence}

Two one-sided tests compare the identity-free tier (T2) against full access (T4). The verdicts depend
on both the model and the metric, and Table~\ref{tab:tost} shows how.

\begin{table}[htbp]
\centering
\caption{TOST equivalence, T2 (linkage only) vs T4 (full surveillance access), $\delta = 0.03$. The
equivalence bound is the smallest margin that would yield \textsc{equivalent} and is margin-free.
$\delta$ was set on the AUC scale; its use on AP is flagged post hoc (Section~\ref{sec:instrument}).}
\label{tab:tost}
\begin{tabular}{llccl}
\toprule
Model & Metric & $\Delta$ (90\% CI) & Equiv.\ bound & Verdict \\
\midrule
Gradient boosting & AUC & $+0.001$ [$-0.000$, $+0.002$] & 0.002 & \textsc{equivalent} \\
                  & AP  & $+0.004$ [$-0.001$, $+0.010$] & 0.010 & \textsc{equivalent} \\
\addlinespace
Logistic & AUC & $+0.011$ [$+0.007$, $+0.015$] & 0.015 & \textsc{equivalent} \\
         & AP  & $+0.062$ [$+0.048$, $+0.078$] & 0.078 & \textsc{surveillance\ superior} \\
\bottomrule
\end{tabular}
\end{table}

For the gradient-boosted model, T2 and T4 are equivalent on both metrics, with an average-precision
equivalence bound of 0.010: full surveillance access buys the capable detector essentially nothing
over pseudonymous linkage, and that holds at any margin above one AP point.

For the logistic model the two metrics part ways, and the disagreement is diagnostic rather than
awkward. On AUC the comparison reads \textsc{equivalent} ($+0.011$); on average precision it reads
\textsc{surveillance superior} ($+0.062$, the entire interval above the margin). The same scores, the
same tiers, and the same test return opposite verdicts because AUC has saturated and average precision
has not. This is the concrete form of the warning in Section~\ref{sec:instrument}: the retracted
analysis certified its central null on AUC, the one metric that cannot see the effect it was used to
rule out.

The intervals in Table~\ref{tab:tost} condition on a single generating-process realization and a
single cross-validation fit; they capture entity-level resampling variance but not the variance from
the generator seed or the fold assignment. The equivalence bound of 0.010 is therefore a
within-one-world quantity. The sweep of the next subsection bounds the missing generator variance
directly: across its fifteen conditions the gradient-boosted T2-vs-T4 average-precision equivalence
bound ranges from 0.010 to 0.036, so the capable model's identity increment stays within four
average-precision points of zero even under the least favorable seed and obfuscation level.

\subsection{Robustness across the generating range}
\label{sec:robustness}

Because the deepest objection to a synthetic ablation is that its answer is an artifact of the chosen
generating process, we do not report the result at a single operating point. We sweep the process's
most consequential free parameter---laundering \texttt{obfuscation}, the degree to which laundering
entities spread and disguise their behavior---across five levels from 0.5 to 0.9, at three seeds each,
holding everything else fixed. Fifteen conditions per model. The degeneracy audit passes at every one.

The point estimates wander, as point estimates from 400-odd positive cases do. The \emph{verdict} does
not. Table~\ref{tab:robustness} reports the T2-vs-T4 equivalence outcome on average precision across
all fifteen conditions.

\begin{table}[htbp]
\centering
\caption{T2-vs-T4 equivalence on average precision across the obfuscation sweep (5 levels $\times$ 3
seeds). The invariant, not the point estimates, is the result.}
\label{tab:robustness}
\begin{tabular}{lccl}
\toprule
Model & \textsc{surv.\ superior} & \textsc{equivalent} / \textsc{inconcl.} & $\Delta$ range \\
\midrule
Logistic          & 15 / 15 & 0 / 15 & $+0.048$ to $+0.076$ \\
Gradient boosting & 0 / 15  & 14 / 1 & $+0.004$ to $+0.025$ \\
\bottomrule
\end{tabular}
\end{table}

Across the whole range, the logistic model finds full surveillance access decisively superior to
linkage---\textsc{surveillance superior} in every one of fifteen conditions---and the gradient-boosted
model never does, in any condition, on either metric (0 of 60). The capable model's average-precision
verdict is \textsc{equivalent} in fourteen of fifteen conditions and \textsc{inconclusive} in one
(highest obfuscation, one seed); it is never \textsc{surveillance superior}. So the strongest honest
statement about the capable model is that the identity increment is equivalent to zero or, at worst,
not distinguishable from it---never that surveillance wins. The gap between the two models is not a
feature of one world. It is an invariant of the generating process across the range we can move it.

\subsection{The instrument can return the opposite answer}
\label{sec:falsify}

An ablation harness that reports ``identity adds nothing'' on every input is measuring nothing. The
falsification world \texttt{surveillance\_strong} is parameterized so that identity attributes
genuinely dominate behavioral structure, and the pipeline aborts unless the harness detects it. It
does: on average precision, T2-vs-T4 returns \textsc{surveillance superior} for both models
($+0.145$ gradient boosting, $+0.437$ logistic). The same code that reports equivalence in the primary
world reports the opposite here.

One detail in this world is worth stating because it closes the argument of
Section~\ref{sec:equivalence} from the other side. For the gradient-boosted model in the falsification
world, average precision returns \textsc{surveillance superior} ($+0.145$) while AUC returns
\textsc{inconclusive} ($+0.033$). Here surveillance genuinely does win, by construction---and AUC
cannot see it. A metric that fails to detect a deliberately planted effect cannot be used to certify
that effect's absence, which is precisely what the original AUC-only analysis did.

\subsection{Interpretation}
\label{sec:interpretation}

The marginal value of identity access is not a fixed property of financial data. It is a function of
the detector applied to that data: a model that cannot reconstruct behavior from structure needs
identity attributes to substitute for what it cannot infer, and a model that can reconstruct behavior
finds the same attributes add nothing measurable. Detection capability and identity collection are to
a meaningful degree substitutes, and a monitoring regime that treats identity access as a fixed
requirement is holding its own analytical capability constant without acknowledging that it has done
so. This holds across the generating range, on the metric that a monitoring regime actually operates
on, and it reverses cleanly when the world is built so that it should.

Two qualifications belong here rather than in a later section. The first is that the linkage step,
which carries the capable model's entire gain, is itself a surveillance capability: resolving
pseudonymous wallets to a common owner is exactly the inference pseudonymity is meant to frustrate.
The result of this paper is that the identity increment \emph{above linkage} is small---a genuine
finding, and a narrower privacy claim than ``monitoring is unnecessary.'' The second is that this is
measured on synthetic data, for a reason Section~\ref{sec:limitations} develops: no public
anti-money-laundering dataset carries an identity axis, so the quantity decomposed here cannot at
present be measured on real data by us or anyone else. The synthetic instrument is not a convenience;
it is currently the only instrument that can pose the question at all.
\section{What the Measurement Licenses, and What It Does Not}
\label{sec:licenses}

The result of Section~\ref{sec:results} is narrow, and it is worth being precise about what
architectural conclusions it supports before the rest of the paper draws them. The measurement
licenses one claim and forbids several tempting extensions of it.

\subsection{What it licenses}

For a detector able to exploit the interaction structure of behavior, full identity access is
statistically equivalent to pseudonymous linkage alone, robustly across the generating range: the
identity increment above linkage is equivalent to zero, or at worst indistinguishable from it, and
never favorable to surveillance (Section~\ref{sec:robustness}). What this licenses is an architecture
in which detection operates on \emph{linked but pseudonymous} activity---an entity's wallets resolved
to a common owner whose civil identity is not attached---and identity is disclosed only on a
warranted, per-case basis rather than collected by default. The detection layer specified in
Section~\ref{sec:architecture} is built to that specification, and the measurement is what
justifies it: identity-by-default is not buying the detection performance it is assumed to buy.

This also answers the objection that graph-only detection is simply worse and the privacy is paid for
in missed crime. At the operating point, the capable detector's identity-free tier reaches an area
under the curve of 0.998 and average precision of 0.980, statistically equivalent to full access; the
performance is not sacrificed, it is recovered from structure. The privacy-preserving configuration
is not a worse detector with a privacy consolation. It is, for a capable model, the same detector.

\subsection{What it does not license}

Three extensions do not follow, and the architecture in the following sections is careful not to
assume them.

First, the result is about the increment \emph{above linkage}, and linkage is itself a surveillance
capability. Resolving an entity's pseudonymous wallets to a common owner is exactly the inference
that unlinkable pseudonymity is designed to frustrate, and it carries most of the capable detector's
performance. The honest reading is not that monitoring is unnecessary but that the \emph{identity}
component of monitoring, above the linkage component, adds little. An architecture that leans on this
result must still take a deliberate position on linkage---who may perform it, under what authority,
and whether it happens by default or on cause---rather than treating the paper as a license to link
freely because ``identity adds nothing.''

Second, the equivalence is a property of a capable detector, and the same measurement shows a
capacity-limited detector obtaining decisive value from identity access at every point of the
generating range. An architecture that withholds identity data on the strength of this result is
implicitly committing its operator to maintaining detection capability at the level the result
assumes. That is a reasonable commitment, but it is a commitment, and it should be stated as an
operational requirement rather than assumed away.

Third, the finding is measured on synthetic data, and the falsification world demonstrates that a
differently constructed process returns the opposite verdict. The result is evidence that identity
adds little \emph{in a regime where laundering is a latent behavior recoverable from structure}; it
is not evidence that this regime describes any particular real deployment. Section~\ref{sec:limitations}
takes up why the real-data question cannot presently be settled and what would settle it.
\section{The Architecture the Measurement Licenses}
\label{sec:architecture}

The measurement supports an architecture in which detection operates on linked-but-pseudonymous
activity and civil identity is disclosed only on cause. This section specifies that architecture. Its
cryptographic components are drawn from the privacy-preserving-compliance lineage of
Section~\ref{sec:already-known} rather than invented here; what the measurement adds is a reason to
prefer this configuration over an identity-by-default one, namely that identity-by-default is not
buying the detection value it is assumed to buy. Two design elements---the abandonment primitive and
opt-in deanonymization---are where the privacy headroom the measurement identifies is actually spent.

\subsection{Design principles}

\textbf{Separation of pattern detection from identity.} Transaction-graph analysis can identify
structural anomalies---velocity, unusual counterparty structure, layering motifs---without knowing
participant identities. The detection system outputs suspicion scores for transactions, never for
users.

\textbf{Transaction-level intervention.} Suspicious activity affects the specific flagged
transaction, not the user's broader financial access. A flagged payment is delayed; the user's
ability to transact otherwise is untouched.

\textbf{Opt-in deanonymization.} Identity revelation is voluntary. A user may explain a flagged
transaction, disclosing whatever they consider appropriate, or abandon it without explanation and
without any record linking identity to the flagged pattern.

\textbf{Architectural enforcement.} Privacy guarantees are structural, not policy. The analytic
components do not possess the capability to link transactions to civil identities without active user
participation, which is a stronger guarantee than a promise to refrain from using such a capability.

\subsection{Detection, intervention, and resolution}

The \textbf{Detection Layer} ingests anonymized transaction data---amounts, timestamps, pseudonymous
counterparty tokens---and maintains a rolling transaction graph it cannot deanonymize. A pattern
engine compares structures against publicly auditable typology libraries, an anomaly system flags
novel structures, and a scorer emits per-transaction suspicion scores. Its output is ``transaction
$X$ has suspicion score $Z$,'' never ``user $Y$ is suspicious.''

The \textbf{Intervention Layer} freezes flagged transactions in escrow and notifies both parties
through the ordinary confirmation channel, without additional identifying information. This
notification is where the architecture departs from the current reporting regime, and
Section~\ref{sec:tipping-off} states that tension rather than eliding it. Parties may explain, request
time, or abandon; none of these requires identity disclosure. On timeout the transaction is canceled,
funds return to sender, and no record links any identity to the flagged pattern.

The \textbf{Resolution Layer} receives voluntary explanations together with a zero-knowledge proof
that the explainer is a party to the transaction, without revealing which party. Trained reviewers see
the explanation, the suspicion rationale, and the amount; they do not see identities, account
histories, or other transactions. Decisions---release or maintain freeze---are logged for audit, and
anonymized outcomes feed back into detection so the system learns from resolved cases without holding
identities.

\subsection{The identity firewall}

The firewall is what makes the privacy guarantee structural. Transaction-graph data and identity
mapping are held by separate systems with no shared access; identity lives only in the user-facing
Wallet Layer and is never transmitted to the analytic components. When a user chooses to explain, a
zero-knowledge proof establishes party status without revealing identity. There is no backdoor: the
components that perform analysis simply do not hold identity information, so state actors cannot
compel per-transaction identity revelation without user cooperation or threshold activation. Where
governance authorizes exceptional access---a court-ordered investigation---a threshold scheme requires
multiple independent parties to authorize decryption for specific transactions, making bulk
surveillance technically rather than merely legally unavailable.

\subsection{The abandonment mechanism and its deterrent}

The abandonment primitive is where the measurement's privacy headroom is spent, and its strategic
structure is worth making explicit. The analysis here is informal: it identifies the choices a flagged
actor faces and reasons about incentives, without deriving equilibria, and Section~\ref{sec:limitations}
takes up the repeated-interaction dynamics it leaves open.

A flagged actor has three options: explain and proceed, which for a genuinely illicit transaction
requires producing a legitimate narrative and may create evidence; abandon silently, canceling the
transaction with funds returned and no identity linked; or wait for timeout, equivalent to
abandonment with delay. The deterrent follows from the first two. A laundering operation that cannot
move money has failed in its core objective, regardless of whether anyone is identified, and unlike
current systems---where sophisticated actors ``comply through'' by producing documentation---there is
no comply-through option here: either a transaction can be legitimately explained, in which case it
may well be legitimate, or it cannot proceed. Aggregate abandonment patterns also feed detection
without identifying anyone.

For legitimate users the friction is low. Detection tuned for high-confidence anomalies flags a small,
governance-controlled fraction of legitimate transactions; a flagged legitimate payment usually has a
ready explanation that need only be plausible, not formally verified; abandonment imposes no penalty
beyond transaction failure for a user who prefers not to explain; and a flag never carries
account-level consequence. Relative to surveillance-based designs, this trades away a class of harms:
no account freezes or investigations from false positives, no repurposing of monitoring infrastructure
for ends beyond its original mandate, and no honeypot of identity-linked transaction data to be
breached, because that data is never assembled.

\subsection{Targeted investigation remains available}

The architecture prevents mass retrospective search, not investigation. Subpoenas can compel
individuals to produce records; court orders can require wallet providers to disclose identity for
specific users; undercover methods remain; and threshold activation can enable identity access for
specific transactions when governance authorizes it. What is foreclosed is bulk access---the ability
to search all transactions retrospectively for patterns associated with individuals---which is not
required for targeted investigation and is the specific capability whose collection the public
objects to. \citet{Rennie2021} distinguish four privacy ``losses'' a CBDC can impose---anonymity,
liberty, individual control, and regulatory control---and argue each needs separate treatment rather
than aggregation; the state's capacities differ qualitatively from commercial data collection, since
commercial entities cannot freeze property or incarcerate \citep{IMF2024CBDC}. The architecture
implements the principle that identity access be gated by due process structurally rather than by
policy commitment.

\subsection{Adoption and scale}

The objection that no state would adopt a privacy-preserving CBDC has weakened. The U.S. Anti-CBDC
Surveillance State Act (H.R.~1919 / S.~1124, 119th Congress), which passed the House in 2025,
prohibits the Federal Reserve from issuing a CBDC usable ``as a tool for surveilling Americans'';
\citet{Zatti2025} reads this as a shift from frameworks asking how CBDCs should be designed to ones
rejecting surveillance architectures outright. The Act does not endorse the architecture proposed
here, but it establishes that state opposition to surveillance CBDC has moved from civil-society
concern onto the legislative agenda. Per Congress.gov the bill has passed the House; House passage
is not enactment, and neither that step nor the introduction of S.~1124 makes either bill law. The paper's claim is in any case normative rather than predictive:
demonstrating that a privacy-preserving alternative exists forecloses the argument that surveillance is
a technical necessity, whatever any given jurisdiction chooses. On scale, the architecture assumes
most transactions proceed without intervention and only flagged ones require attention; at a flag rate
below a fraction of a percent, a system clearing millions of transactions daily produces thousands of
flags, within reach of automated triage with human review reserved for contested cases. High-frequency
or institutional flows can be placed under a distinct regime---accepting identity registration in
exchange for reduced friction---without compromising retail privacy.

\subsection{An unresolved tension: tipping-off}
\label{sec:tipping-off}

One objection has force the architecture does not dissolve. FATF Recommendation~21 requires that
institutions and their staff be prohibited by law from disclosing that a suspicious-transaction report
is being filed \citep{FATF2012}, and national implementations such as section~333A of the UK Proceeds
of Crime Act~2002 criminalize disclosures likely to prejudice an investigation. The anonymous
notification at the center of the Intervention Layer does what these provisions forbid: it tells both
parties, in real time, that their transaction has drawn suspicion. The honest statement is that the
architecture is incompatible with the current reporting regime, not that it satisfies it. It
presupposes a different regime in which prevention substitutes for covert reporting---the flagged
transaction cannot complete, so there is no downstream investigation for notification to prejudice in
the usual sense---and whether regulators would accept that substitution is an open legal question that
would require amending FATF standards before any deployment. Notification also interacts with the
probing problem of Section~\ref{sec:limitations}: telling parties which transactions are flagged is
exactly what makes the detection boundary observable. A deployable design might delay or coarsen
notification, but each such change erodes the due-process property that motivates it. We record this
as an unresolved compliance-design tension rather than a solved problem.
\section{Feasibility of the Stack}
\label{sec:feasibility}

The measurement of Section~\ref{sec:results} licenses an architecture (Section~\ref{sec:architecture})
in which detection runs on linked-but-pseudonymous activity and identity is disclosed only on cause.
This section addresses a separate question: whether that architecture is buildable from primitives
that exist today. The answer we defend is feasibility, not optimality, and the distinction matters
for how the simulation at the end of the section should be read---it establishes that the
\emph{composition} of known components stays within a plausible performance envelope, not that any
component's cost has been independently validated.

\subsection{Cryptographic requirements}

The architecture relies on established primitives. Zero-knowledge proofs provide transaction
validation and party authentication; the Zcash deployment demonstrates production viability of
comparable requirements \citep{Sasson2014}. Secure multi-party computation enables pattern detection
without any single party holding complete data. Threshold cryptography gates emergency identity
access behind multiple independent authorizations. Blind signatures enable authorization without
linkage; \citet{ChaumGrothoffMoser2021} give a complete CBDC issuance protocol in which the central
bank cannot link withdrawals to spending. Ring signatures allow a party to prove membership in a
transaction set without revealing which member it is, useful for anonymous explanation in the
Resolution Layer.

\subsection{Feasibility assessment}

Each primitive has been demonstrated in production or at research scale: zero-knowledge proofs in
Zcash's shielded transaction pool \citep{ElectricCoin2023}; secure multi-party computation in
privacy-preserving analytics such as Google's Private Join and Compute \citep{Google2019}; threshold
cryptography in cryptocurrency custody and distributed key management \citep{Gennaro2016}; blind
signatures across privacy-preserving payment systems since their introduction \citep{Chaum1982}.

Performance remains a genuine consideration. Current ZKP systems carry significant proving overhead
relative to transparent transactions, though proof generation can run client-side and verification is
far cheaper than generation. Project Hamilton, the Federal Reserve Bank of Boston and MIT
collaboration, showed a CBDC transaction processor sustaining on the order of $10^6$ transactions per
second with sub-second latency \citep{Lovejoy2023Hamilton}, which establishes headroom for the
additional cost of privacy-preserving primitives at national scale even though that throughput figure
was measured without those primitives in the path.

A parallel ZKP-based compliance design merits comparison. \citet{decker2025zkp} propose a
zero-knowledge KYC architecture that reduces exposed user data substantially while reporting high
fraud-detection accuracy, in the narrower setting of point-of-entry institutional verification. Their
setting and ours are complementary rather than competing: their construction could serve as the
identity-verification layer at account creation, while the architecture here governs ongoing
transaction monitoring without identity exposure. We do not compare detection rates across the two,
because they measure different quantities on different populations.

Progress toward practical viability continues. \citet{OfflineCBDC2024} present a Secure Element-based
offline CBDC system with latency comparable to commercial payment rails, and advances in secure
multi-party computation continue to reduce communication overhead \citep{MPC2024}. Zero-knowledge
tooling in particular has matured into a substantial engineering base across blockchain platforms
\citep{ElHajj2024ZKBench}, though claims about universal industry adoption should be read with the
usual caution about vendor reporting.

\subsection{System architecture}

A complete system separates five layers, each with a defined interface and no shared identity access
beyond the first. The \textbf{Wallet Layer} is the only component holding user identity; it generates
proofs and interfaces with the layers below. The \textbf{Transaction Layer} processes anonymized
transactions and settles the ledger. The \textbf{Detection Layer} analyzes patterns on pseudonyms.
The \textbf{Intervention Layer} manages flagged transactions and resolution, with identity access
only where a user voluntarily discloses. The \textbf{Governance Layer} manages parameters, threshold
key shares, and audit, exercising oversight without operational access. Compromise of any single
layer does not enable identity compromise absent user cooperation or threshold activation.

\subsection{A concrete PET--AML stack}
\label{sec:pet-aml-stack}

The remaining subsections instantiate the high-level claim into a privacy-enhancing AML (PET--AML)
stack that can be built today: private watchlist checking via Private Set Intersection (PSI) or
VOPRF-based membership tests; secure risk propagation over pseudonymous transaction graphs; and
transaction-time policy proofs in zero-knowledge. These map onto the layered architecture---the
Wallet Layer screens and proves, the Transaction Layer verifies, the Detection Layer scores
pseudonyms, and the Governance Layer controls exceptional exposure.

\subsubsection{Threat model}

We assume a wallet/PSP performing KYC and issuing an unlinkable credential; a ledger operator
validating transaction proofs and maintaining settlement; a compliance authority maintaining
watchlists and policy parameters; and a governance quorum holding threshold key shares for
exceptional identity exposure. The adversaries of concern are honest-but-curious operators, malicious
users seeking to launder while remaining unlinkable, and partial compromise or collusion among
operational entities. There is deliberately no single super-admin trust anchor: identity access
requires threshold activation and legal authorization. The security objectives are transaction
privacy (the ledger validates payments and enforces limits without learning identity), watchlist
privacy (neither the compliance authority nor the ledger learns non-matching queries or the list
itself), enforceability (no valid transaction without satisfied policy constraints), and due process
(disclosure is auditable, rate-limited, threshold-gated, with no bulk de-anonymization interface).

\subsubsection{Components}

\paragraph{Private watchlist checking (PSI/VOPRF).} Sanctions and PEP screening are implemented as a
privacy-preserving membership test. The wallet derives a stable identifier from KYC material and runs
a PSI-style protocol against the watchlist. Unbalanced PSI is well-suited to this setting, in which
the server holds a large set and each client query set is small \citep{Wang2025UnbalancedPSI}; an
alternative uses verifiable OPRFs \citep{rfc9497}, with the wallet checking membership against a
periodically published PRF-transformed watchlist. On a no-match the wallet can carry a compact
witness proving screening occurred without revealing the identifier; on a match, policy triggers a
dispute window in the Intervention Layer rather than automatic identity leakage.

\paragraph{Secure risk propagation on pseudonymous graphs.} The Detection Layer benefits from
graph-structured signals but should not learn identities, so identity is treated as a sealed label
and risk is computed on pseudonyms. Because even which pseudonyms are seeded as ``known bad'' can be
sensitive, the seed vector is secret-shared while non-identifying graph structure remains available,
and a bounded-iteration risk function is evaluated under secure computation. Cross-institution graphs
can use secure graph-analytics frameworks that reduce edge-proportional communication overhead and
have been demonstrated at million-node scale under semi-honest models \citep{Koti2024Graphiti,
Yu2025GraphAce}. A privacy-preserving MPC pilot among Dutch banks reported by \citet{vanEgmond2024}
gives some evidence that computing risk signals across institutional boundaries, without sharing
customer identities, is practically approachable rather than only theoretical.

\paragraph{ZK policy proofs at admission.} Transaction-time proofs make policy constraints verifiable
on entry: balance correctness and no double-spend; per-tier limits on amount and cumulative spend via
range and sum proofs; credential validity with selective disclosure, optionally using threshold-issued
credentials such as Coconut \citep{Sonnino2019Coconut}; and proof that screening was performed for the
current epoch. SNARKs offer small proofs and fast verification at the cost of trusted setup, while
STARKs are transparent and post-quantum secure with larger proofs \citep{ElHajj2024ZKBench}; a
jurisdiction can select a point on that frontier and upgrade over time.

\subsubsection{Performance envelope}

Table~\ref{tab:pet-aml-envelope} gives an order-of-magnitude envelope. It is indicative, not
normative, and the figures are drawn from the cited benchmarks rather than measured here.

\begin{table}[H]
\centering
\caption{Indicative performance envelope for the PET--AML stack (order-of-magnitude; figures from the
cited sources, not measured in this work).}
\label{tab:pet-aml-envelope}
\begin{tabular}{p{3.2cm} p{3.2cm} p{3.2cm} p{4.0cm}}
\toprule
\textbf{Operation} & \textbf{Latency} & \textbf{Bandwidth / storage} & \textbf{Notes} \\
\midrule
Watchlist check (wallet $\rightarrow$ FIU) & sub-second interactive; amortizable via epochs & tens of KB per check; server stores $|W|$ items & Unbalanced PSI supports large $|W|$ at modest client cost \citep{Wang2025UnbalancedPSI} \\
\addlinespace
Transaction ZK proof generation & sub-second to a few seconds on commodity hardware & proof size $\sim\!10^2$ bytes (SNARK) to $\sim\!10^4$ bytes (STARK) & Benchmarked range spans roughly 0.5\,ms--3.6\,s across circuits; verification typically milliseconds \citep{ElHajj2024ZKBench} \\
\addlinespace
Secure risk propagation (batch) & minutes to a few hours per batch (e.g.\ daily) & state proportional to $|V|$; preprocessing amortizes multiplications & Secure graph frameworks report million-scale feasibility and large comms reductions \citep{Koti2024Graphiti,Yu2025GraphAce} \\
\bottomrule
\end{tabular}
\end{table}

To keep routine payments smooth, screening runs at wallet ``ready'' time, proofs are generated at
transaction time from fixed audited circuits, and risk propagation runs in batches unless a high-risk
typology requires near-real-time scoring.

\subsubsection{Simulation: feasibility of the composition}
\label{sec:pet-aml-simulation}

To check that the composition of these components stays within the envelope, we implement a
discrete-event simulation of the stack (\texttt{pet\_aml\_sim.py}).\footnote{Source in the paper's
repository, \url{https://github.com/andrewmaksakov/CBDC}.} The simulator models 40{,}000 wallets
(eight PSPs of 5{,}000) across four KYC tiers in proportions $55/30/13/2\%$, draws transaction amounts
from a lognormal distribution, and exercises PSI screening, ZK proof verification, and batch MPC risk
propagation. Table~\ref{tab:sim-results} reports a representative run.

\begin{table}[H]
\centering
\caption{PET--AML simulation, representative run (40{,}000 transactions, seed 7).}
\label{tab:sim-results}
\begin{tabular}{p{6.4cm} p{3.4cm}}
\toprule
\textbf{Metric} & \textbf{Value} \\
\midrule
Transactions generated       & 40{,}000 \\
Transactions settled         & 25{,}287 (63.2\%) \\
Transactions rejected        & 14{,}713 (36.8\%) \\
\quad rejected under tier-conditioned amounts (sensitivity) & 3{,}375 (8.4\%) \\
Sanctions hits (PSI blocked) & 82 \\
Escalations (case packets)   & 27 (0.068\% of screened; 14 later settled) \\
\addlinespace
Latency mean / p50 / p90 / p99 & 537 / 511 / 768 / 1{,}137\,ms \\
Latency breakdown            & PSI screening dominates \\
MPC batch runtime            & 0.38\,hours \\
\bottomrule
\end{tabular}
\end{table}

\textbf{What the simulation shows, and what it does not.} It shows that the composition is feasible:
given per-component costs taken from the literature, queueing, contention, batching, and cross-PSP
routing do not push end-to-end latency outside the retail envelope---mean 537\,ms, 99th percentile
1.14\,s, dominated by PSI screening. This is a statement about the architecture's composition, not a
validation of the component costs themselves, which enter the simulator as inputs
(Table~\ref{tab:pet-aml-envelope}); a simulation parameterized by an envelope cannot confirm that
envelope, and we do not claim it does.

The 36.8\% rejection rate is an artifact of the amount distribution, not a system failure mode:
14{,}631 of the 14{,}713 rejections are tier-0 per-transaction-limit breaches (\$50 cap) and the daily
cap never binds, because amounts are drawn from one global lognormal whose median exceeds that cap
irrespective of tier. Drawing amounts per-wallet from a tier-scaled distribution (mean 40\% of the
wallet's tier limit, \texttt{--tier-amounts}) drops the rate to 8.4\%, with per-transaction latency
essentially unchanged; the rejection rate is therefore set by the ratio of typical spend to the cap
rather than by the architecture. The 82 sanctions blocks confirm that PSI screening intercepts
watchlist matches without revealing match status to the ledger.

The escalation count requires care in interpretation. An earlier version of the simulator reported
zero escalations, an artifact of processing order: risk tiers were assigned during batch propagation
\emph{after} the transaction loop, so tier-based escalation could never fire. With propagation moved
to run once per simulated day, the same two-day run yields 27 escalations. These are counted at the
point of triggering, before settlement, so they are not a subset of the settled transactions---14 of
the 27 subsequently settle and 13 are rejected, which is why the rate is reported against screened
rather than settled transactions. The figure is a small, real exercise of the escalation pathway, not
evidence about escalation rates in a deployed system, which would depend on the true prevalence and
behavior of high-risk wallets.
\section{Governance}
\label{sec:governance}

Privacy-preserving does not mean ungoverned. Several parameters require ongoing oversight rather than
technical default. Detection thresholds trade higher sensitivity against more false positives, a
policy choice with distributional consequences. Timeout periods trade accommodation of legitimate
users against friction on illicit abandonment. Pattern libraries determine what triggers a flag and
should be publicly auditable to permit scrutiny and prevent discriminatory targeting. And the
holders of threshold key shares---who can authorize exceptional identity access---should be
distributed across independent institutions so that capture is hard and emergency capability is
preserved.

Transparency is what makes this oversight real. Aggregate statistics on flag rates, resolution
outcomes, and abandonment patterns can be published without compromising any individual, and they
are what let the public assess whether the system works. Resolution decisions should be periodically
reviewed by independent auditors for consistency and abuse; system source code should be public and,
where feasible, formally verified; and every activation of exceptional identity access should be
logged with justification, with unauthorized attempts triggering investigation.

These are not merely technical settings. \citet{KeisterSanches2023} show through general-equilibrium
analysis that CBDC design choices carry non-trivial consequences for banking stability, deposit
funding, and credit allocation, so the architectural properties of digital currency are determinants
of macroeconomic structure rather than implementation details. When choices carry consequences of
that magnitude they belong to democratic processes rather than technical defaults. The framework is
consistent with stakes-weighted accounts of legitimacy, in which governance should track stakeholder
preferences in proportion to the stakes involved \citep{farzulla2025consensual}: financial-privacy
architecture bears more acutely on marginalized communities, dissidents, and the economically
vulnerable than on those with institutional access, and governance should reflect that heterogeneity.
The evolving European framework illustrates both the difficulty and the feasibility of legislating
privacy requirements, as the digital euro's offline mode raises particular AML questions under the
Anti-Money-Laundering Regulation, MiCAR, and PSD3 that architectural privacy guarantees could
simplify relative to layered exemptions \citep{Minto2024}. Governance structures should accordingly
include legislative authorization specifying privacy requirements, an independent oversight body with
audit authority, public reporting on operation and decisions, mechanisms for citizen input on
parameters, and sunset clauses requiring periodic reauthorization.
\section{Limitations and Future Work}
\label{sec:limitations}

The central limitation is structural rather than incidental, and it bounds what any version of this
study can claim. The quantity we decompose---the marginal value of identity access, separated from
the linkage beneath it---requires a dataset in which entities carry both a linkage structure and
verified identity attributes, with laundering labels. No public anti-money-laundering dataset has
this shape. The two standard benchmarks, the IBM AMLworld synthetic transaction data and the
Elliptic Bitcoin graph, both lack an identity axis entirely: they support the linkage question
(T1$\to$T2) and not the identity question (T2$\to$T3$\to$T4). The consequence is unavoidable and we
state it plainly: the paper's central quantity cannot at present be measured on real data by us or by
anyone else, because the data that would measure it does not exist in the public domain. This is why
a controlled synthetic instrument is not a convenience but a necessity---it is the only setting in
which the identity axis can be varied at all.

That said, the synthetic setting is the study's second limitation, and the honest response to it is
not realism but generator independence. The result currently rests on a data-generating process of
our own construction, and the degeneracy audit, the falsification world, and the obfuscation sweep
are all defenses against the objection that we designed the process to produce the answer. The
strongest remaining defense is to run the parts of the decomposition that public data can support on
public data. Future work will apply the linkage tier (T1$\to$T2) of the harness to AMLworld HI-Small,
whose account structure suits an account-based CBDC, marking T3 and T4 not-evaluable and reporting the
linkage result on an externally generated benchmark. What this buys is not realism---AMLworld is
itself synthetic---but independence from our generator on the one axis a public dataset can exercise,
which is the direct answer to the objection that sank the earlier version of this work.

The third limitation is statistical. The canonical operating point realizes 388 positive cases, which
is enough to identify the tier increments and to render the equivalence verdicts stable across seeds,
but it is not large, and the average-precision intervals remain wide enough that the study leans on
the equivalence test rather than on differencing point estimates. A larger positive class would
tighten the surface further; it would not change the invariant, which is already stable across the
sweep, but it would narrow the one inconclusive cell.

Finally, the architecture the measurement licenses carries a limitation of its own that the
measurement does not touch. The abandonment primitive---allowing a user to walk away from a flagged
transaction without deanonymization---doubles as a free adversarial probing oracle: an actor can
submit transactions to learn the detection boundary and abandon those that are flagged, refining
evasive behavior at no cost. This is a genuine tension between the privacy guarantee and the
detection guarantee, it is not resolved here, and Section~\ref{sec:architecture} notes it at the point
where the primitive is specified rather than deferring it to a caveat.


\section{Conclusion}
\label{sec:conclusion}

Proposals for retail central bank digital currency have largely accepted that anti-money-laundering
obligations require identity-linked transaction monitoring, and have taken the design problem to be
how much privacy can be preserved around that requirement. This paper asked what the requirement is
worth, and measured it. Decomposing surveillance access into three increments---pseudonymous linkage,
identity attributes, and watchlist consultation---and pricing each against a fixed detector on fixed
data, we find that the marginal value of identity access is not a property of the data but of the
detector: a model that can reconstruct behavior from structure recovers almost everything from
linkage and gains nothing measurable from identity, while a model that cannot leans on identity to
substitute for what it cannot infer. Detection capability and identity collection are, to a
meaningful degree, substitutes.

The finding is deliberately hedged where the evidence requires it. It holds robustly across the range
of the generating process's most consequential parameter, on the metric a monitoring regime actually
operates on, and it reverses cleanly when the process is rebuilt so that identity genuinely dominates.
It is confined to the increment above linkage, and linkage is itself a surveillance capability. It is
established on synthetic data, for the structural reason that no public dataset carries the identity
axis the question requires. Within those bounds, the conclusion is that identity collection by default
is not buying the detection value it is assumed to buy, and that whether to collect identity by
default is therefore a policy choice conditioned on detection capability rather than a technical
necessity.

The architecture and cryptographic stack that this measurement licenses, set out in
Sections~\ref{sec:architecture}--\ref{sec:feasibility}, are assembled from a mature
privacy-preserving-compliance literature rather than invented here, and their feasibility rests on
existing primitives. Their contribution is that the measurement now gives a reason to prefer them:
not that surveillance is impossible to build, but that, for a capable detector, the identity-linked
version of it adds little the pseudonymous version does not already achieve. The instrument that
establishes this, and the harness's insistence on auditing itself for degeneracy, testing equivalence
rather than inferring it, and proving it can return the opposite verdict, are offered as much as the
contribution as the number they produce---because in this corner of the literature, the way a result
is manufactured has too often been the thing that made it wrong.
\section*{Declarations}

\paragraph{Conflict of Interest.} The authors declare no competing interests.

\paragraph{Funding.} This research received no external funding.

\paragraph{Data Availability.} All code and results are available at
\url{https://github.com/andrewmaksakov/CBDC}. The detection-ablation harness of
Sections~\ref{sec:instrument}--\ref{sec:results} (\texttt{detection/}) is seeded and self-contained:
\texttt{run\_all.py} regenerates the committed results, \texttt{sweep\_obfuscation.py} regenerates the
robustness surface of Section~\ref{sec:robustness}, and \texttt{requirements.txt} pins the environment
under which the committed figures reproduce. The PET--AML stack simulation of
Section~\ref{sec:feasibility} (\texttt{pet\_aml\_sim.py}) is in the same repository.

\paragraph{AI Assistance.} Claude (Anthropic) was used as a research collaborator for analysis,
literature synthesis, LaTeX preparation, and iterative refinement. All intellectual claims and errors
remain the authors' responsibility.

\clearpage

\bibliography{references}

\end{document}
